% The EI-reflection of a category, its preorder reflection, and the homomorphism order of graphs.
\documentclass[11pt]{amsart}
\usepackage{amsmath,amssymb}
\setlength{\emergencystretch}{3em}
\tolerance=600
\usepackage{tikz-cd}
\usepackage[hidelinks]{hyperref}

\newtheorem{thm}{Theorem}[section]
\newtheorem{lem}[thm]{Lemma}
\newtheorem{cor}[thm]{Corollary}
\newtheorem{prop}[thm]{Proposition}
\theoremstyle{definition}
\newtheorem{defi}[thm]{Definition}
\newtheorem{rem}[thm]{Remark}
\newtheorem{ex}[thm]{Example}
\newtheorem{ques}[thm]{Question}

\newcommand{\Same}{\mathrm{Same}}
\newcommand{\Mor}{\mathrm{Mor}}
\newcommand{\Iso}{\mathrm{Iso}}
\newcommand{\Aut}{\mathrm{Aut}}
\newcommand{\Hom}{\mathrm{Hom}}
\newcommand{\End}{\mathrm{End}}
\newcommand{\Ob}{\mathrm{Ob}}
\newcommand{\Loc}{\mathrm{Loc}}
\newcommand{\Sat}{\mathrm{Sat}}
\newcommand{\supp}{\mathrm{supp}}
\newcommand{\Wcyc}{W_{\mathrm{cyc}}}
\newcommand{\Wobj}{W_{\mathrm{obj}}}
\newcommand{\Wend}{W_{\mathrm{end}}}
\newcommand{\Wret}{W_{\mathrm{ret}}}
\newcommand{\Wsupp}{W_{\mathrm{supp}}}
\newcommand{\Cat}{\mathbf{Cat}}
\newcommand{\Set}{\mathbf{Set}}
\newcommand{\sSet}{\mathbf{sSet}}
\newcommand{\Presh}{\mathrm{Presh}}
\newcommand{\Sub}{\mathrm{Sub}}

\begin{document}

\title[The EI-reflection of a category]{The EI-reflection of a category,
  its preorder reflection, and the homomorphism order of graphs}
\author{Ryota Shibusawa}
\thanks{Corresponding author. E-mail: \texttt{r-shibusawa@daiichi-koudai.ac.jp}}
\address{Daiichi Institute of Technology}
\email{r-shibusawa@daiichi-koudai.ac.jp}
\keywords{antisymmetry, preorder reflection, EI-category, localization,
  two-out-of-three, relative category, subterminal object, homomorphism
  order, core of a graph}
\subjclass[2020]{18A32 (primary), 18A20, 06A75, 18B25, 05C60 (secondary)}

\begin{abstract}
Antisymmetry of an order --- $x\le y$ and $y\le x$ imply $x=y$ --- becomes,
when equality is replaced by a class $W$ of morphisms declared to be
``sameness'', the condition that every morphism admitting a morphism back
lies in $W$. The morphisms admitting a morphism back form a class $\Wcyc(C)$
of any category $C$, closed under composition and two-out-of-three; it is
the largest such class inverted by the preorder reflection, it is $\{\mathrm{id}\}$
exactly for acyclic categories and $\Iso(C)$ exactly for EI-categories, and
its localization $C\to C[\Wcyc^{-1}]$ is the reflection of $C$ into
EI-categories. The main theorem is that when $C$ has binary coproducts (or
binary products) this localization is a preorder: the EI-reflection of $C$
is its preorder reflection $X\preceq Y\iff\Hom(X,Y)\neq\varnothing$. For the
topos of graphs this preorder is the homomorphism preorder of Hell and
Ne\v set\v ril, and on finite graphs the resulting identification is the
passage to cores. We contrast the localization with the reflection: in any
topos the objects local for $\Wcyc$ are the subterminal ones, so the
saturation of $\Wcyc$ is the class of maps with equal support, i.e.\ the
$(-1)$-truncation; $\Wcyc$ is saturated for sets and simplicial sets, where
it is the $(-1)$-truncated equivalences, but not for graphs. Finally the
classes $\Wcyc$, $\Wobj$ (source and target isomorphic), $\Wend$ (identities
and endomorphisms), $\Aut$, $\Iso$ are shown to be \emph{natural} --- preserved
by all functors --- and every natural such class other than the class of all
morphisms is contained in $\Wcyc$. All proofs are elementary and no result
depends on a computation.
\end{abstract}

\maketitle

\section{Introduction}

An ordered set is a preorder in which mutual comparability forces equality.
Replacing that equality by an isomorphism gives the familiar notion of a
skeletal category; replacing it by an identification with more structure
leads to univalent categories and complete Segal spaces. This paper studies
the intermediate, and elementary, situation in which the equality is
replaced by an arbitrary class $W$ of morphisms declared to be a
\emph{sameness}: a class containing the identities, closed under
composition and satisfying two-out-of-three. Such classes form a complete
lattice $\Same(C)$ \cite{Same}, and the question ``for which $W$ is $C$
antisymmetric up to $W$?'' has a canonical answer.

\begin{defi}\label{def:cyc}
For a category $C$ write $X\preceq Y$ if $\Hom(X,Y)\neq\varnothing$ (the
\emph{reachability preorder}). A morphism $f\colon X\to Y$ \emph{has a back}
if $Y\preceq X$. Let
\[
\Wcyc(C):=\{f\colon X\to Y\ :\ \Hom(Y,X)\neq\varnothing\}
\]
be the class of morphisms with a back (equivalently, the morphisms inside
the strongly connected classes of $\preceq$). For $W\in\Same(C)$, the
category $C$ is \emph{$W$-antisymmetric} if $\Wcyc(C)\subseteq W$.
\end{defi}

Thus a preorder is $\{\mathrm{id}\}$-antisymmetric iff it is a partial order,
and a category is $\Iso$-antisymmetric iff mutually related objects are
isomorphic through the relating morphisms. The results are as follows.

\medskip\noindent\textbf{The cycle class and the EI-reflection.} $\Wcyc(C)$
is a sameness, and it is the largest sameness inverted by the preorder
reflection $C\to(\Ob C,\preceq)$ (Theorem~\ref{thm:cyc}). It is $\{\mathrm{id}\}$
iff $C$ is acyclic, $\Iso(C)$ iff $C$ is an EI-category (all endomorphisms
invertible), and $\Mor(C)$ iff $C$ is a coproduct of strongly connected
categories (Proposition~\ref{prop:extremes}). The first new result is that
the localization $C\to C[\Wcyc^{-1}]$ is an EI-category and is the
\emph{reflection} of $C$ into EI-categories (Theorem~\ref{thm:EI}): inverting
every morphism that has a morphism back produces the universal EI-category
under $C$. It is standard that in an EI-category the reachability
preorder on isomorphism classes of objects is a partial order
\cite[\S9]{Lueck}; Theorem~\ref{thm:EI} goes in the other direction, from an
arbitrary category to an EI-category, and identifies the class of
morphisms that must be inverted.

\medskip\noindent\textbf{Main theorem.} The second new result is that if
$C$ has binary coproducts, or binary products, then $C[\Wcyc^{-1}]$ is a
preorder, namely $(\Ob C,\preceq)$ (Theorem~\ref{thm:main}). So for such categories the EI-reflection and the
preorder reflection coincide, an EI-category with binary coproducts is a
preorder, and the antisymmetrization of the topos of directed graphs is the
homomorphism order of graphs \cite{HN}: identifying graphs that map to each
other is the same as passing to cores (Corollary~\ref{cor:graphs}).

\medskip\noindent\textbf{Localization versus reflection.} In a topos the
objects local for $\Wcyc$ are exactly the subterminal ones, so the
saturation of $\Wcyc$ is the class of maps with equal support, the class
inverted by the $(-1)$-truncation (Theorem~\ref{thm:topos}). Hence
``antisymmetry up to $\Wcyc$'' has two shadows: the reflection sees only
supports, the localization sees reachability. They agree, and $\Wcyc$ is
saturated, for sets and for simplicial sets --- where $\Wcyc$ is the class of
$(-1)$-truncated equivalences of \cite{Postnikov} --- but not for graphs.

\medskip\noindent\textbf{Natural samenesses.} The third new result
concerns samenesses defined uniformly for all categories. The classes $\{\mathrm{id}\}$,
$\Aut$, $\Iso$, $\Wend$ (identities and endomorphisms), $\Wobj$ (source
isomorphic to target), $\Wret$ (source and target mutual retracts), $\Wcyc$,
$\Mor$ are samenesses preserved by every functor; they form a lattice, in
which $\Iso$ and $\Wend$ are incomparable with join $\Wobj$; and every
natural sameness other than $\Mor$ is contained in $\Wcyc$
(Theorem~\ref{thm:cut}), so that $\Wcyc$ is the largest proper natural
sameness: a single constraint on every two-out-of-three class that is
defined naturally in the category. No result of this paper depends on a
computation; small cases were explored by computer only to find the
statements.

\medskip\noindent\textbf{Related work.} EI-categories appear in
transformation groups \cite{Lueck} and in the representation theory of
categories \cite{Webb}; we have not seen the observation that they form a
reflective subcategory of $\Cat$ with the localization at $\Wcyc$ as unit.
The preorder and partial-order reflections of a category are standard
constructions; the point here is the intermediate EI-reflection, the class
of morphisms it inverts, and the fact that under (co)products the two
reflections coincide.
Categories with a chosen class of ``weak equivalences'' are the relative
categories of Barwick and Kan \cite{BK} and the homotopical categories of
\cite{DHKS}; samenesses are relative categories whose class satisfies
two-out-of-three, and the lattice $\Same(C)$ is the fibre over $C$ of the
forgetful functor from relative categories to categories
(Remark~\ref{rem:relcat}). The homomorphism order of graphs and the theory
of cores are classical \cite{HN}. Subterminal objects, supports and the
$(-1)$-truncation are standard topos theory \cite{Elephant,MLM}. The
companion papers \cite{Same,Postnikov,Topjoin} introduce the sameness
lattice and study its sheaf-theoretic and homotopical elements; the present
paper is self-contained.

\section{Preliminaries}\label{sec:prelim}

\subsection{Conventions}
Categories are small unless said otherwise (a topos is a large category, and
for it the statements below are applied objectwise or in a larger universe;
nothing depends on this). $\Cat$ denotes the category of small categories.

\subsection{Samenesses}
A \emph{sameness} on a category $C$ is a class $W$ of morphisms with (i)
$\mathrm{id}_X\in W$ for all $X$, (ii) $g\circ f\in W$ whenever $f,g\in W$ are
composable, (iii) two-out-of-three: for composable $f,g$, if two of
$f,g,g\circ f$ lie in $W$ then so does the third. $\Same(C)$ is the set of
samenesses ordered by inclusion; it is closed under arbitrary
intersections and contains $\Mor(C)$, hence is a complete lattice with
$\bot=\{\mathrm{id}\}$ and $\top=\Mor(C)$.

\begin{lem}\label{lem:pullback}
For a functor $F\colon C\to D$ and $V\in\Same(D)$, $F^*(V):=\{f:F(f)\in V\}$ is a
sameness on $C$. In particular $F^*(\Iso(D))$ is the largest sameness
inverted by $F$.
\end{lem}
\begin{proof}
$F$ preserves identities and composites, so the three conditions for $V$
transfer to $F^*(V)$. Any class inverted by $F$ is contained in
$F^*(\Iso(D))$, which is itself inverted.
\end{proof}

\subsection{Localization}
For $W\subseteq\Mor(C)$ the localization $\gamma\colon C\to C[W^{-1}]$ is the
initial functor inverting $W$ \cite[I.1]{GZ}. Its morphisms are represented
by zigzags $X\to A_1\leftarrow A_2\to\cdots\to Y$ in which the backward
arrows lie in $W$, modulo the relations generated by composition and
$w\circ w^{-1}=\mathrm{id}$; in particular $\Hom_{C[W^{-1}]}(X,Y)\neq\varnothing$ iff
such a zigzag exists. Note that $\gamma$ is bijective on objects and that
$C[W^{-1}]$ is small when $C$ is.

\subsection{Local objects and saturation}
For $W\in\Same(C)$ an object $Z$ is \emph{$W$-local} if $\Hom(f,Z)\colon
\Hom(Y,Z)\to\Hom(X,Z)$ is bijective for every $f\colon X\to Y$ in $W$;
$\Loc(W)$ denotes the class of $W$-local objects and
$\Sat(W):=\{f:\Hom(f,Z)\text{ bijective for all }Z\in\Loc(W)\}$ the
\emph{saturation} of $W$, a sameness containing $W$. If $L\colon C\to C$ is a
reflector onto a full replete subcategory $\mathcal E$ then the class
$W_L=\{f:L(f)\text{ iso}\}$ is saturated with $\Loc(W_L)=\mathcal E$: a map is
inverted by $L$ iff $\Hom(f,Z)$ is bijective for all $Z\in\mathcal E$, and a
$W_L$-local $Z$ has $\Hom(\eta_Z,Z)$ bijective, whence $\eta_Z\colon Z\to LZ$ is
a split monomorphism with a retraction $r$, and $\eta_Zr=\mathrm{id}_{LZ}$ because
$\Hom(\eta_Z,LZ)$ is injective; so $Z\cong LZ\in\mathcal E$.

\begin{rem}[relative categories]\label{rem:relcat}
A pair $(C,W)$ with $W$ a wide subcategory is a relative category
\cite{BK}; requiring two-out-of-three gives a full subcategory
$\mathbf{RelCat}_{2/3}$. The forgetful functor $\mathbf{RelCat}_{2/3}\to\Cat$ is a
bifibration whose fibre over $C$ is $\Same(C)$, with cartesian lifts $F^*$
and cocartesian lifts $F_!$ (the sameness generated by the image); the
functor $\Same\colon\Cat\to\mathbf{CLat}_{\mathrm{Gal}}$ of \cite{Same} is its
Grothendieck construction. Nothing below depends on this remark.
\end{rem}

\section{The cycle class}\label{sec:cyc}

\begin{thm}\label{thm:cyc}
For every category $C$, $\Wcyc(C)$ is a sameness, and
$\Wcyc(C)=P^*(\Iso)$ for the preorder reflection $P\colon C\to(\Ob C,\preceq)$.
Hence $\Wcyc(C)$ is the largest sameness inverted by the preorder (and by
the partial-order) reflection of $C$, and
\[
\{\mathrm{id}\}\ \subseteq\ \Iso(C)\ \subseteq\ \Wobj(C)\ \subseteq\ \Wcyc(C)\ \subseteq\ \Mor(C),
\]
where $\Wobj(C):=\{f\colon X\to Y: X\cong Y\}$.
\end{thm}
\begin{proof}
A morphism of the preorder $(\Ob C,\preceq)$ is invertible iff its endpoints
are mutually related, so $P(f)$ is an isomorphism iff $f$ has a back:
$\Wcyc(C)=P^*(\Iso)$, a sameness by Lemma~\ref{lem:pullback}, and the largest
one inverted by $P$. The partial-order reflection is $P$ followed by the
skeleton, which inverts the same morphisms. Identities and isomorphisms have
backs; if $X\cong Y$ then any $f\colon X\to Y$ has a back; so the chain holds
($\Wobj$ is a sameness by Lemma~\ref{lem:objdet}).
\end{proof}

Directly: if $f\colon X\to Y$ and $g\colon Y\to Z$ have backs $h,k$ then $gf$ has
the back $hk$; if $gf$ has a back $m$ and $f$ has a back $h$ then $g$ has the
back $fm$; if $gf$ has a back $m$ and $g$ has a back $k$ then $f$ has the
back $mg$. Figure~\ref{fig:cyc} shows the three cases.

\begin{figure}[ht]
\centering
\begin{tikzcd}[column sep=large]
X\arrow[r,"f"] & Y\arrow[r,"g"]\arrow[l,bend left=35,"h"] & Z\arrow[l,bend left=35,"k"]
\end{tikzcd}
\qquad
\begin{tikzcd}[column sep=large]
X\arrow[r,"f"] & Y\arrow[r,"g"]\arrow[l,bend left=35,"h"] & Z\arrow[ll,bend left=50,"m"]
\end{tikzcd}
\caption{Two-out-of-three for $\Wcyc$: backs compose (left); if $gf$ and $f$
have backs $m$ and $h$, then $f\circ m$ is a back for $g$ (right); the third
case is symmetric.}
\label{fig:cyc}
\par\smallskip\noindent\textit{Alt text:} Two diagrams of three objects in a
row with forward arrows $f$ and $g$; curved arrows return from $Y$ to $X$ and
from $Z$ to $Y$ in the first, and from $Y$ to $X$ and from $Z$ to $X$ in the
second.
\end{figure}

\begin{prop}[the three extremes]\label{prop:extremes}
For a category $C$:
\begin{enumerate}
\item $\Wcyc(C)=\{\mathrm{id}\}$ iff $C$ is \emph{acyclic}: $X\preceq Y\preceq X$
implies $X=Y$, and the only endomorphisms are identities;
\item $\Wcyc(C)=\Iso(C)$ iff $C$ is an \emph{EI-category}: every endomorphism
is invertible;
\item $\Wcyc(C)=\Mor(C)$ iff $C$ is a coproduct of strongly connected
categories, i.e.\ $\preceq$ is symmetric.
\end{enumerate}
\end{prop}
\begin{proof}
(1) and (3) are restatements. (2) If $C$ is EI and $f\colon X\to Y$ has a back
$g$, then $gf\in\End(X)$ and $fg\in\End(Y)$ are invertible, so $f$ has the left
inverse $(gf)^{-1}g$ and the right inverse $g(fg)^{-1}$, which then coincide:
$f$ is invertible. Conversely every endomorphism has the identity as a back,
so $\Wcyc(C)=\Iso(C)$ forces endomorphisms to be invertible.
\end{proof}

\begin{thm}[EI-reflection]\label{thm:EI}
For every small category $C$ the localization $C[\Wcyc^{-1}]$ is an
EI-category with the same reachability preorder as $C$, and
$\gamma\colon C\to C[\Wcyc^{-1}]$ is universal among functors from $C$ to
EI-categories. Hence the full subcategory of (small) EI-categories is reflective in
$\Cat$, with unit the localization at $\Wcyc$.
\end{thm}
\begin{proof}
\emph{Convexity.} If $X\preceq A\preceq Y$ and $Y\preceq X$ then $A$ is in the
$\preceq$-class of $X$. Consequently every zigzag in $C$ between two objects
of one class, with backward arrows in $\Wcyc$, stays inside that class: a
backward arrow $A_{i+1}\to A_i$ in $\Wcyc$ has a back, so $A_i\preceq A_{i+1}$
as well, and inductively all $A_i$ lie between $X$ and $Y$.

\emph{Reachability.} A zigzag from $X$ to $Y$ with backward arrows in $\Wcyc$
yields $X\preceq Y$ by replacing each backward arrow by a back; conversely a
morphism $X\to Y$ is a zigzag. So $\Hom_{C[\Wcyc^{-1}]}(X,Y)\neq\varnothing$ iff
$X\preceq Y$.

\emph{EI.} An endomorphism of $X$ in $C[\Wcyc^{-1}]$ is represented by a
zigzag from $X$ to $X$, which by convexity lies in the class of $X$; every
$C$-morphism between objects of that class has a back, hence is inverted by
$\gamma$; so the zigzag is a composite of isomorphisms and their inverses,
and is invertible.

\emph{Universality.} Let $G\colon C\to E$ with $E$ an EI-category. If
$f\colon X\to Y$ has a back $g$, then $G(f)$ has the back $G(g)$, so $G(f)$ is
invertible by Proposition~\ref{prop:extremes}(2). Thus $G$ inverts $\Wcyc$ and
factors uniquely through $\gamma$.
\end{proof}

\begin{rem}
For $W\supseteq\Wcyc$ the further localization $C[W^{-1}]$ need not be an
EI-category: in the poset with elements $a<b$, $a<d$, $c<b$, $c<d$ and
$W=\{a<b,\ c<d\}$ (so $W\supseteq\Wcyc=\{\mathrm{id}\}$), the zigzag
$a\to d\leftarrow c\to b\leftarrow a$ is an endomorphism of $a$ in
$C[W^{-1}]$ which is not invertible, as the localization of this poset at
the two edges is the free category on a loop at $a$ together with the
remaining arrows. $W$-antisymmetry is a property of the pair $(C,W)$, not of
the localization.
\end{rem}

\section{Main theorem: categories with binary coproducts}\label{sec:main}

\begin{thm}\label{thm:main}
Let $C$ be a category with binary coproducts, or with binary products.
Then $C[\Wcyc^{-1}]$ is a preorder: it is the reachability preorder
$(\Ob C,\preceq)$. Consequently the EI-reflection of $C$ is its preorder
reflection, and the skeleton of $C[\Wcyc^{-1}]$ is the partial-order
reflection of $C$.
\end{thm}
\begin{proof}
Let $g,g'\colon B\to C'$ be parallel morphisms of $C$, and let
$j_1,j_2\colon B\to B\sqcup B$ be the coproduct injections and
$\nabla\colon B\sqcup B\to B$ the fold map, so that
$\nabla j_1=\nabla j_2=\mathrm{id}_B$. The morphism $\nabla$ has the back $j_1$, so
$\nabla\in\Wcyc$ is inverted in $C[\Wcyc^{-1}]$, where therefore
$j_1=\nabla^{-1}=j_2$. Hence
\[
g=[g,g']\circ j_1=[g,g']\circ j_2=g'\qquad\text{in }C[\Wcyc^{-1}],
\]
where $[g,g']\colon B\sqcup B\to C'$ is the copairing. So any two parallel
morphisms of $C$ become equal, and then any two parallel zigzags become
equal as well (a zigzag is a composite of morphisms of $C$ and inverses of
morphisms of $C$, and inverses of equal morphisms are equal); thus
$C[\Wcyc^{-1}]$ has at most one morphism between any two objects. By
Theorem~\ref{thm:EI} it has a morphism $X\to Y$ iff $X\preceq Y$. With
binary products use instead the projections $p_1,p_2\colon C'\times C'\to C'$,
which have the common back $\Delta$, so $p_1=\Delta^{-1}=p_2$ and
$g=p_1\langle g,g'\rangle=p_2\langle g,g'\rangle=g'$
(Figure~\ref{fig:main}).
\end{proof}

\begin{figure}[ht]
\centering
\begin{tikzcd}[column sep=large,row sep=large]
B\arrow[r,shift left,"j_1"]\arrow[r,shift right,"j_2"'] & B\sqcup B\arrow[l,bend right=45,"\nabla"']\arrow[r,"{[g,g']}"] & C'
\end{tikzcd}
\caption{Theorem~\ref{thm:main}: $\nabla$ has the back $j_1$ and is inverted;
both injections then equal $\nabla^{-1}$, and $g=[g,g']j_1$, $g'=[g,g']j_2$
coincide.}
\label{fig:main}
\par\smallskip\noindent\textit{Alt text:} An object $B$ with two parallel
arrows into the coproduct of $B$ with itself, a curved fold arrow back to
$B$, and the copairing arrow from the coproduct to $C'$.
\end{figure}

\begin{ex}[the hypothesis is needed]
Let $C$ be the free category on the graph $x\rightleftarrows y$. Every
non-identity morphism has a back, so $\Wcyc=\Mor(C)$ and $C[\Wcyc^{-1}]$ is the
free groupoid on the graph, whose endomorphism group at $x$ is $\mathbb Z$.
It is an EI-category but not a preorder.
\end{ex}

\begin{cor}\label{cor:EIco}
An EI-category with binary coproducts (or binary products) is a preorder;
in particular a groupoid with binary coproducts is equivalent to a discrete
category, and a monoid with binary coproducts is trivial.
\end{cor}
\begin{proof}
In an EI-category the fold map $\nabla$, having the back $j_1$, is already
invertible in $C$ (Proposition~\ref{prop:extremes}), so the computation in
the proof of Theorem~\ref{thm:main} takes place in $C$: $j_1=\nabla^{-1}=j_2$
and $g=g'$ for all parallel $g,g'$. A groupoid and a monoid are
EI-categories.
\end{proof}

\begin{cor}[toposes and presheaf categories]\label{cor:topos}
For a topos $\mathcal E$, and in particular for every presheaf category
$\Presh(S)$, the EI-reflection $\mathcal E[\Wcyc^{-1}]$ is the hom-preorder
of $\mathcal E$: $X\preceq Y$ iff there is a morphism $X\to Y$.
\end{cor}

\begin{cor}[graphs and cores]\label{cor:graphs}
Let $\mathbf{Gph}$ be the category of directed graphs (presheaves on the
category with two objects and two parallel non-identity arrows). Then
$\mathbf{Gph}[\Wcyc^{-1}]$ is the homomorphism preorder of graphs, $X\preceq Y$
iff there is a graph homomorphism $X\to Y$; its partial-order reflection is
the homomorphism order of \cite{HN}. Two finite graphs are identified in
$\mathbf{Gph}[\Wcyc^{-1}]$ iff they have isomorphic cores.
\end{cor}
\begin{proof}
The first statement is Corollary~\ref{cor:topos}. Finite graphs $X,Y$ are
mutually reachable iff they are homomorphically equivalent, iff their
cores are isomorphic \cite[Ch.~1]{HN}.
\end{proof}

So, for finite graphs, ``antisymmetry up to isomorphism'' is exactly the
passage from a graph to its core, and the sameness lattice records this
passage as the single element $\Wcyc\in\Same(\mathbf{Gph})$.

\section{Natural samenesses}\label{sec:natural}

\begin{defi}
A \emph{natural sameness} is an assignment $C\mapsto W(C)\in\Same(C)$, for all
small categories $C$, such that $F(W(C))\subseteq W(D)$ for every functor
$F\colon C\to D$; equivalently, a subfunctor of the functor
$\Mor\colon\Cat\to\Set$ that is pointwise a sameness.
\end{defi}

Natural samenesses are closed under pointwise intersection, and under
pointwise join in $\Same(C)$ (the image of a join is contained in the join
of the images, because the sameness generated by $F(W\cup W')$ contains
$F(W\vee W')$: the class $F^*(W(D)\vee W'(D))$ is a sameness containing $W(C)$
and $W'(C)$). So they form a lattice $\mathbf{NatSame}$ mapping into every
$\Same(C)$ by lattice homomorphisms.

\begin{lem}\label{lem:objdet}
Let $R$ be an equivalence relation on the objects of $C$. Then
$W_R:=\{f\colon X\to Y: X\,R\,Y\}$ is a sameness (\emph{object-determined}); if
$R$ is preserved by all functors then $W_R$ is natural. Equality gives
$\bot$, isomorphism gives $\Wobj$, mutual reachability gives $\Wcyc$.
\end{lem}
\begin{proof}
Reflexivity gives the identities, transitivity gives closure under
composition, and transitivity with symmetry gives two-out-of-three
($X\,R\,Z$ and $X\,R\,Y$ imply $Y\,R\,Z$, and similarly in the other cases).
\end{proof}

\begin{prop}\label{prop:naturals}
The following are natural samenesses:
\begin{itemize}
\item $\bot=\{\mathrm{id}\}$, $\Iso$, $\Wobj$ (source isomorphic to target), $\Wcyc$,
$\top=\Mor$;
\item $\Wend(C):=\{\mathrm{id}\}\cup\{\text{endomorphisms of }C\}$ and
$\Aut(C):=\{\mathrm{id}\}\cup\{\text{automorphisms}\}=\Iso\cap\Wend$;
\item $\Wret(C):=\{f\colon X\to Y: X\text{ and }Y\text{ are retracts of each other}\}$.
\end{itemize}
Moreover $\Iso$ and $\Wend$ are incomparable in general, $\Iso\vee\Wend=\Wobj$,
$\Wobj\subseteq\Wret\subseteq\Wcyc$, and $\Wret=\Wobj$ in every category whose
endomorphism monoids are finite.
\end{prop}
\begin{proof}
Functors preserve identities, isomorphisms, isomorphic pairs of objects,
backs, endomorphisms, automorphisms and retractions, so all classes are
natural once they are samenesses. $\Wobj$, $\Wret$: Lemma~\ref{lem:objdet},
``mutual retract'' being an equivalence relation containing $\cong$ (a retract
of a retract is a retract). $\Wend$: composites of endomorphisms of $X$ are
endomorphisms; if $gf$ and $f$ are endomorphisms of $X$ then $g\colon X\to X$;
if $gf$ and $g$ are endomorphisms then $\mathrm{dom}\,f=\mathrm{cod}\,gf=\mathrm{cod}\,g
=\mathrm{dom}\,g=\mathrm{cod}\,f$. $\Aut$ is an intersection. Incomparability: in
the category with two objects and one isomorphism between them,
$\Iso\not\subseteq\Wend$; in the monoid $\{1,t,t^2\}$ with $t^3=t^2$,
$\Wend\not\subseteq\Iso$. Join: $\Wobj$ is a sameness containing both; and if
$f\colon X\to Y$ with an isomorphism $h\colon Y\to X$, then $f=h^{-1}\circ(hf)$ with
$hf\in\Wend$, $h^{-1}\in\Iso$, so $f$ lies in the join. If $X$ is a retract of
$Y$ via $s\colon X\to Y$, $r\colon Y\to X$, $rs=\mathrm{id}$, then $e\mapsto ser$ is an
injection $\End(X)\to\End(Y)$ (with left inverse $u\mapsto rus$); with finite
endomorphism monoids, mutual retracts give a bijection, so
$\mathrm{id}_Y=ser$ for some $e$, and $s$ is a split epimorphism as well as a
split monomorphism, hence an isomorphism.
\end{proof}

\begin{figure}[ht]
\centering
\begin{tikzcd}[row sep=small,column sep=small]
& \top=\Mor & \\
& \Wcyc\arrow[u,dash] & \\
& \Wret\arrow[u,dash] & \\
& \Wobj\arrow[u,dash] & \\
\Iso\arrow[ur,dash] & & \Wend\arrow[ul,dash] \\
& \Aut\arrow[ul,dash]\arrow[ur,dash] & \\
& \bot\arrow[u,dash] &
\end{tikzcd}
\caption{Natural samenesses of Proposition~\ref{prop:naturals}; $\Iso$ and
$\Wend$ are incomparable with meet $\Aut$ and join $\Wobj$. Further natural
samenesses exist between $\Wobj$ and $\Wcyc$ (Remark~\ref{rem:classify}).}
\label{fig:nat}
\par\smallskip\noindent\textit{Alt text:} A Hasse diagram: bottom, then
automorphisms, then two incomparable elements (isomorphisms; endomorphisms),
their join (source isomorphic to target), then mutual retracts, then the
cycle class, then the top.
\end{figure}

\begin{thm}[the cut lemma]\label{thm:cut}
Let $W$ be a natural sameness. If for some category $C$ the class $W(C)$
contains a morphism $f\colon X\to Y$ with $Y\not\preceq X$, then $W=\top$.
Hence every natural sameness other than $\Mor$ is contained in $\Wcyc$, and
$\Wcyc$ is the largest proper natural sameness. Dually, if $W(C)$ contains a
morphism $f\colon X\to Y$ with $X\neq Y$, then $W\supseteq\Iso$.
\end{thm}
\begin{proof}
Let $U:=\{Z: Z\preceq X\}$ and $V:=\Ob C\setminus U$; then $X\in U$, $Y\in V$, and
no morphism goes from $V$ to $U$ (if $Z'\to Z$ with $Z\preceq X$ then
$Z'\preceq X$). So $U\mapsto 0$, $V\mapsto1$ defines a functor $F\colon C\to[1]$
to the arrow category, with $F(f)$ the non-identity arrow $u$ of $[1]$.
Naturality gives $u\in W([1])$, and every morphism of every category is the
image of $u$ under a functor $[1]\to D$; so $W=\top$. For the second
statement let $I$ be the category with two objects $0,1$ and exactly one
morphism between any two objects (the walking isomorphism). Any assignment
of objects of $C$ to $\{0,1\}$ extends uniquely to a functor $C\to I$; choose
$X\mapsto0$, $Y\mapsto1$, so $F(f)$ is the isomorphism $u\colon0\to1$ and
$u\in W(I)$. An isomorphism $h\colon A\to B$ of any $D$ is the image of $u$
under the functor $I\to D$ sending $0\mapsto A$, $1\mapsto B$, $u\mapsto h$,
$u^{-1}\mapsto h^{-1}$ (also when $A=B$). So $W\supseteq\Iso$.
\end{proof}

\begin{rem}[towards a classification]\label{rem:classify}
By Lemma~\ref{lem:objdet}, natural equivalence relations on objects give
natural samenesses $W_R$; by the cut lemma every such relation other than the
total one is contained in mutual reachability. There are infinitely
many: the transitive closure of any natural symmetric relation defined by
a finite pattern (mutual retracts; pairs $u\colon X\to Y$, $v\colon Y\to X$ with
$uvu=u$ and $vuv=v$; \dots) is one. Not every natural sameness is of this
form ($\Wend$ is not), and two-out-of-three couples the object part with
the endomorphism part ($\Iso\vee\Wend=\Wobj$). We do not know the complete list. Restricting a natural sameness
$W\subseteq\Wend$ to one-object categories yields a natural assignment
$M\mapsto N(M)$ of two-out-of-three-closed submonoids of monoids, of which we
know only $\{1\}$, the group of units, and $M$; whether this restriction
determines $W$ --- i.e.\ whether membership of an endomorphism in $W(C)$
depends only on its image in $\End_C(X)$ --- is a further classification
question.
\end{rem}

\section{Localization versus reflection in a topos}\label{sec:topos}

Let $\mathcal E$ be a topos. Every object $X$ has a \emph{support}
$\supp X\rightarrowtail1$, the image of $X\to1$; the subterminal objects (subobjects
of $1$) form a reflective full subcategory $\Sub(1)$ with reflector
$X\mapsto\supp X$, the $(-1)$-truncation \cite[A1.3]{Elephant}. Put
$\Wsupp:=\{f\colon X\to Y:\supp X=\supp Y\}$, the class inverted by this
reflector.

\begin{thm}\label{thm:topos}
In a topos $\mathcal E$:
\begin{enumerate}
\item $\Loc(\Wcyc)=\Sub(1)$, the subterminal objects;
\item $\Sat(\Wcyc)=\Wsupp$;
\item $\Wcyc\subseteq\Wsupp$, with equality iff any two objects with the
same support admit morphisms in both directions.
\end{enumerate}
\end{thm}
\begin{proof}
(1) Let $Z$ be subterminal. For any $Y$, $\Hom(Y,Z)$ is empty or a singleton,
and it is non-empty iff $\supp Y\le Z$ in $\Sub(1)$ (a map $Y\to Z$ factors
through the image of $Y\to1$, and conversely). If $f\colon X\to Y$ has a back
$g$, then $\supp X\le\supp Y$ because of $f$ and $\supp Y\le\supp X$ because of
$g$; so $\Hom(Y,Z)=\Hom(X,Z)$ and $Z$ is $\Wcyc$-local. Conversely let $Z$ be
$\Wcyc$-local and $X$ arbitrary. The fold map $\nabla\colon X\sqcup X\to X$
has the back $j_1$, so $\nabla\in\Wcyc$ and $\Hom(X,Z)\cong\Hom(X\sqcup X,Z)
=\Hom(X,Z)\times\Hom(X,Z)$ via the diagonal; a set in bijection with its
square through the diagonal has at most one element. So $\Hom(X,Z)$ has at
most one element for every $X$, i.e.\ $Z\to1$ is a monomorphism: $Z$ is
subterminal.
(2) $\Wsupp$ is the class inverted by the reflector onto $\Sub(1)$, hence is
saturated with local objects $\Sub(1)$; by (1) it has the same local objects
as $\Wcyc$, so it is $\Sat(\Wcyc)$.
(3) The inclusion was shown in (1). If $f\colon X\to Y$ lies in $\Wsupp$ and
objects with equal support admit maps both ways, then $f$ has a back; the
converse is the definition.
\end{proof}

\begin{figure}[ht]
\centering
\begin{tikzcd}[column sep=large,row sep=large]
\mathcal E\arrow[r,"\gamma"]\arrow[d,"\supp"'] & \mathcal E[\Wcyc^{-1}]=(\Ob\mathcal E,\preceq)\arrow[dl,dashed] \\
\Sub(1) &
\end{tikzcd}
\caption{Localization versus reflection at $\Wcyc$ in a topos: the
localization is the hom-preorder (Corollary~\ref{cor:topos}); the reflection
onto the local objects is the support, i.e.\ the $(-1)$-truncation
(Theorem~\ref{thm:topos}); the dashed comparison is an equivalence iff
$\Wcyc$ is saturated.}
\label{fig:topos}
\par\smallskip\noindent\textit{Alt text:} A triangle: the topos maps to its
hom-preorder by the localization and to the subterminal objects by the
support; a dashed arrow compares the hom-preorder with the subterminals.
\end{figure}

\begin{ex}[sets and simplicial sets]\label{ex:sset}
In $\Set$ and in $\sSet$ the supports are $\varnothing$ and $1$, and two
non-empty objects admit maps in both directions (in $\sSet$ through a
vertex, $Y\to\Delta^0\to X$). So $\Wcyc=\Wsupp$ is saturated, it is the class
of maps $f\colon X\to Y$ with ($X=\varnothing$ iff $Y=\varnothing$), i.e.\ the
$(-1)$-truncated equivalences, and $\mathcal E[\Wcyc^{-1}]\simeq\Sub(1)=\{\varnothing<1\}$.
For simplicial sets this class is the join of the weak homotopy
equivalences with the $0$-coskeletal equivalences \cite[Thm.~6.2, $n=-1$]{Postnikov}:
the antisymmetry sameness of simplicial sets is the $(-1)$-truncation.
\end{ex}

\begin{ex}[graphs: $\Wcyc$ is not saturated]\label{ex:graphs}
In $\mathbf{Gph}$ let $X$ be the path $v_0\to v_1$ and $Y$ the graph with one
vertex and one loop (the terminal graph). Both have support $1$, and there
is a map $X\to Y$; but $\Hom(Y,X)=\varnothing$, since a loop must map to a
loop. So $X\to Y$ lies in $\Wsupp\setminus\Wcyc$: $\Wcyc$ is strictly smaller
than its saturation, the localization $\mathbf{Gph}[\Wcyc^{-1}]$ (the
homomorphism preorder, Corollary~\ref{cor:graphs}) is strictly finer than
the reflection $\mathbf{Gph}\to\Sub(1)$, and the comparison of
Figure~\ref{fig:topos} identifies all graphs with at least one edge.
\end{ex}

The pattern is the one observed in \cite{Topjoin} for joins of Grothendieck
topologies: the sameness lattice sees a relation --- here reachability, there
a pro-limit --- that the saturated, reflective level forgets.

\section{Concluding remarks}

The cycle class $\Wcyc$ answers the question ``antisymmetric up to which
sameness?'' with a canonical element of $\Same(C)$, whose extremes classify
acyclic, EI and strongly connected categories, whose localization is the
EI-reflection, and which under binary coproducts is the preorder reflection
in disguise. Two questions remain. First, the classification of natural
samenesses (Remark~\ref{rem:classify}): is every natural sameness generated
by natural equivalence relations on objects and natural
two-out-of-three-closed submonoids of endomorphism monoids? Second, for
which categories without binary coproducts is the EI-reflection still a
preorder --- the free category on $x\rightleftarrows y$ shows that some
hypothesis is needed, and finite categories with a zero object satisfy the
conclusion by the same fold-map argument applied to the zero morphisms.

\section*{Funding}
No funding was received for this work.

\section*{Declarations}
\noindent\textbf{Competing interests.}\quad The author declares that he has no
competing interests.

\begin{thebibliography}{9}
\bibitem{BK} C.~Barwick, D.~M.~Kan, \emph{Relative categories: another model
for the homotopy theory of homotopy theories}, Indag. Math. (N.S.)
\textbf{23} (2012), 42--68.
\bibitem{DHKS} W.~G.~Dwyer, P.~S.~Hirschhorn, D.~M.~Kan, J.~H.~Smith,
\emph{Homotopy Limit Functors on Model Categories and Homotopical
Categories}, Math. Surveys Monogr. \textbf{113}, Amer. Math. Soc., 2004.
\bibitem{GZ} P.~Gabriel, M.~Zisman, \emph{Calculus of Fractions and Homotopy
Theory}, Ergebnisse der Mathematik \textbf{35}, Springer, 1967.
\bibitem{HN} P.~Hell, J.~Ne\v set\v ril, \emph{Graphs and Homomorphisms},
Oxford Lecture Series in Mathematics and its Applications \textbf{28},
Oxford University Press, 2004.
\bibitem{Elephant} P.~T.~Johnstone, \emph{Sketches of an Elephant: A Topos
Theory Compendium}, Oxford Logic Guides \textbf{43--44}, Oxford University
Press, 2002.
\bibitem{Lueck} W.~L\"uck, \emph{Transformation Groups and Algebraic
$K$-Theory}, Lecture Notes in Mathematics \textbf{1408}, Springer, 1989.
\bibitem{MLM} S.~Mac Lane, I.~Moerdijk, \emph{Sheaves in Geometry and Logic},
Universitext, Springer, 1992.
\bibitem{Same} R.~Shibusawa, \emph{The sameness lattice of a category:
weak-equivalence structures, Galois functoriality, and comparison intervals},
preprint (2026), archived at DOI 10.5281/zenodo.22823772.
\bibitem{Postnikov} R.~Shibusawa, \emph{Postnikov truncations as joins of
sheaf and homotopy equivalences}, preprint (2026), archived at DOI
10.5281/zenodo.22956181.
\bibitem{Topjoin} R.~Shibusawa, \emph{Meets and joins of Grothendieck
topologies in the sameness lattice}, preprint (2026), archived at DOI
10.5281/zenodo.22978446.
\bibitem{Webb} P.~Webb, \emph{An introduction to the representations and
cohomology of categories}, in: Group Representation Theory, EPFL Press,
2007, 149--173.
\end{thebibliography}
\end{document}
